\section{Lösungsansatz}\label{sec:loesungsansatz}
Dieser Abschnitt beschreibt einen der größeren Teile der Arbeit. Hier soll detailliert gezeigt werden, wie das System
zur automatischen Abisolierung in der Prototypenproduktion aussehen soll und welche Entscheidungen auf diesem Weg
getroffen wurden.\\
Dafür sollten noch einmal grob die Anforderungen und Bedingungen dargelegt werden:\\
Das System, das entwickelt werden soll, muss in der Lage sein, aus einer hängenden Reihe von Hairpins einzelne dieser
an den Roboter zu übergeben. Dabei dürfen die Hairpins nicht beschädigt werden und die Übergabe an den Roboterarm
muss mit einer festen Taktzeit und einer hohen Wiederholgenauigkeit geschehen.\\
Zunächst musste sich ein grober Aufbau einer solchen Anlage angesehen werden. Dafür wurden zwei primäre zu erledigende
Arbeitsschritte festgelegt:\\
1. Die Vereinzelung der Hairpins,\\
2. Die Auslieferung und Übergabe eines einzelnen Hairpins an den Roboterarm.\\
Der genaue Aufbau dieser Subsysteme wird im späteren Teil dieses Abschnittes weiter erläutert. Zunächst muss man sich
einen Überblick über die technischen Möglichkeiten verschaffen. Es wird in drei grobe Kategorien unterteilt:\\
- die mechanischen Bauteile \\
- die elektronischen Komponenten \\
- die Programmierung/der Code \\


Es gibt drei Themenbereiche der Anlage:
1. Die Vereinzelung der Hairpins,
2. Die Auslieferung des Hairpins zum Roboter,
3. Das Greifen und Abisolieren des Hairpins mittels Roboterarm, wobei diese Aufgabe nicht Teil dieser Arbeit ist.

\subsection{Mechanische Bauteile}\label{sec:Mechanik}
In diesem Teil werden mechanische Bauteile, die für ein solches System benötigt werden, untersucht. Relevant
sind dabei besonders die Punkte wie Einsetzbarkeit, Langlebigkeit und Praktikabilität.
%
\subsubsection{Aktuatorik}\label{Aktuatorik}

%
\subparagraph{Pneumatik}
%
\subparagraph{Elektronische Aktuatorik}
-Servomotor
-Spindel
-Belt driven
-Roboterarm?
%
%
\subsubsection{Material}\label{sec:MaterialAuswahl}
PLA
%
\subsubsection{Aluminiumprofile}\label{sec:AluProfil}
%
%
\subsection{Elektronische Komponenten}\label{sec:ElektronischeKomponenten}
%
\subsubsection{Spannungsversorgung}\label{sec:SpannungsVersorgung}
%
\subsubsection{Microcontroller}\label{sec:Microcontroller}
%
\subsubsection{Sensorik}\label{sec:Sensorik}
%
\subsubsection{Motoren}\label{sec:Motoren}
%


\subsection{CodeStruktur}\label{sec:CodeStruktur}
Welche Code-Strukturen gibt es für Ausfallsicherheit etc.? Welche Code-Struktur wird verwendet? Welche Vorteile und Nachteile hat diese Code-Struktur?

\subsubsection{Aktuelle Systeme aus der Produktion}\label{sec:Codesystems}
Welche Normen und Arten für solche Systeme gibt es?

\subsubsection{Code-Struktur der Anlage}\label{sec:CodeStrukturAnlage}
Diagramm für Signalfluss, Fallbacks, Sicherheit


\subsection{Vereinzelung}\label{sec:Vereinzelung}
Das Subsystem der Vereinzelung muss die Hairpins mit einer hohen Taktzeit
und Zuverlässigkeit vereinzeln. Dabei muss aus einem Bündel von Hairpins ein
einzelner Hairpin ausgegeben werden.
%
\subsection{Auslieferung}\label{sec:Auslieferung}
Das Subsystem der Auslieferung muss die Hairpins mit einer hohen Taktzeit
und Zuverlässigkeit ausliefern.